% Nouveau EGA V (cote 25): archive cote 25, site page 47 (PDF page 48).
% typescript. Draft compared: none.
% AI-transcribed by Claude Opus 5.5 (claude-opus-5-5, adaptive thinking, effort high), Message Batches API, batch ega25b (msgbatch_01A7n5YuXubUPNK1zsepi8RG), 2026-11-13; not yet checked by eye.
\scanpage{25/47}
\margin{(\unsure{$\uparrow$|?} la variété \unsure{lin} L de codim 2 des hyperplans tangents à X en x, donc par raisons de dim. on a T = L, donc $X \subset L^{\circ}$ = polaire de L, qui est une droite, donc X est un \unsure{\uline{ouvert}} d'une droite.)}
\mspage{48}
\unsure{\hand{(}}En langage classique, elle s'exprime \xout{xxxdisant} sauf erreur en disant
que X est \uline{réglée} pour l'immersion projective considérée \hand{[}- et si ça nous
chante, nous avons ici  tout ce qu'il faut, grâce à 8.5. et ses corollai
res, pour expliquer et justitffier une telle terminologie, au cas où cela
t'inspirerait de faire le lien avec la taupe\over{)}{]}. \inserted{Par exemple, si dim $X=1$, cela implique que X est une droite, car si $x \in X(k)$, alors T contient \unsure{$\uparrow$|P}}

b)  Lorsque la car. est p >0, il faudrait pour bien faire donner
des exemples (avec dim P = 2, X une courbe non singulière) où  les
conditions de 8.6. ne sont pas vérifiées i.e. $\underline{Y}^{\mathrm{sing}} = \underline{Y}^{\mathrm{sup\ sing}}$, et où
néanmoins  dim T = r-1 , i.e. \inserted{des exemples} où  L/K est une extension finie
inséparable.  J'ai la flemme de construire des exemples, mais je ne
doute guère qu'il y en ait.

\deleted{\uline{Corollaire} 8.11.} Dans a), faire un appel du pied au \xout{xxxxx} nº suivant,
où on montrera que lorsque le cas exceptionnel ``réglé'' se produit, alors \margin{N.B. X est dit ``réglé'' pour f si $\mathrm{Im}(\underline{Y}^{\mathrm{sing}} \to \check{P})$ est de codim $\unsure{\geq} 2$}
par une modification triviale de l'immersion projective considérée, on
se retrouve dans la situation  ``générale''  de 8.8.

\margin{\unsure{$\|$}}La partie du présent nº de 8.\over{\illegible{1}}{6}. à ici  devrait sans doute faire
un nº à part, de nature différentielle,  tandis que le début du nº avec
ce qui va suivre serait bloqué dans un nº  sur  la dimension \unsure{dus|des}
exceptionnels. \hand{(Je n'utilise que le fait que $\underline{Y}^{\mathrm{sing}}$ est de
dim r-1).}

\uline{Proposition} 8.1\over{2}{3}.   Nous supposons tou\xout{x}jours f:X$\to$P non ramifié, \inserted{et X sans pt isolé,} \deleted{Suppo-}
\inserted{On suppose que X satisfait $(R_k)$ géométrique\deleted{, et si}}
\deleted{sons que X satisfasse la propriété $(R_k)$ géométrique  (où k \unsure{$\geq$} 0)\hand{,} et si}
\deleted{X' est \xout{xxensemble} l'ouvert de X des points où X est lisse sur k, \unsure{le}}
\deleted{\xout{\illegible{12}}}
\deleted{\inserted{on suppose que le}morphisme  induit f':X'  P  \xout{soit} satisfasse \xout{auxxconditionsxéquivalentes}}
\deleted{\xout{\illegible{10}} k(  )}
\deleted{\xout{dex8x8xxx(disonsxdoncxsimplementxX'xestxnonxx}  à la condition}
\deleted{g( $\underline{Y}'^{\,\mathrm{\unsure{sup\,sing}}}$) \hand{\unsure{´}}de codimension $\geq$2 dans $\hand{$\check{P}$}$ (condition \uline{automatiquement}}
\deleted{\xout{v}vérifiée si k est de car 0 en vertu de de 8.11. ).}  Soit $Z_k$ la parti
de \hand{$\check{P}$} \xout{fermédesxxxxxx} complémentaire de l'ensemble des \hand{$\xi$}$\in$\hand{$\check{P}$}  tels

\margin{47}
